\documentclass[11pt, a4paper]{article}

% --- Packages ---
\usepackage[utf8]{inputenc}
\usepackage{geometry}
\geometry{a4paper, margin=1in}
\usepackage{graphicx} % For including figures (PDF/SVG recommended)
\usepackage{amsmath}  % For math equations
\usepackage{hyperref} % For clickable links to code/references
\usepackage{listings} % For appending code snippets if needed
\usepackage{float}
\usepackage{xcolor}

% --- Setup for Code Listings ---
\lstset{
    basicstyle=\ttfamily\small,
    breaklines=true,
    backgroundcolor=\color{lightgray!10},
    frame=single,
    numbers=left,
    numberstyle=\tiny
}

% --- Document Info ---
\title{Digital Signal Processing: Mandatory Project \\ \large Analysis of the 15 January 2022 Hunga Tonga Eruption}
\author{Your Name / Candidate Number}
\date{September 24, 2026}

\begin{document}

\maketitle

\tableofcontents
\newpage

% ---------------------------------------------------------
\section{Task 1: The Overview}
% ---------------------------------------------------------

\subsection{Dataset Parsing and Station Map (Task 1a)}
% Run script, parse dataset, include map plot
The provided dataset containing the HDF5 sensor time series was parsed. Figure \ref{fig:station_map} displays the geographic map of the event location and the respective recording stations. 

\begin{figure}[H]
    \centering
    % \includegraphics[width=0.8\textwidth]{figures/your_map_plot.pdf} 
    \caption{Geographic locations of the sensors and the Hunga Tonga volcano.}
    \label{fig:station_map}
\end{figure}

\subsection{Shortest and Longest Great Circle Distances (Task 1b)}
Based on the station metadata, the calculated great circle distances from the Hunga Tonga event are as follows:
\begin{itemize}
    \item \textbf{Shortest distance:} [Insert distance here] km (Station: [Insert Station Name/ID])
    \item \textbf{Longest distance:} [Insert distance here] km (Station: [Insert Station Name/ID])
\end{itemize}

\subsection{Distance Distribution Plot (Task 1c)}
Figure \ref{fig:sorted_distances} displays all sorted distances between the event location and the 201 station locations.

\begin{figure}[H]
    \centering
    % \includegraphics[width=0.8\textwidth]{figures/sorted_distances.pdf}
    \caption{Sorted great circle distances from Hunga Tonga to all recording stations.}
    \label{fig:sorted_distances}
\end{figure}

% ---------------------------------------------------------
\section{Task 2: Data Filtering}
% ---------------------------------------------------------

\subsection{FIR Filter Impulse Responses (Task 2a)}
The three provided FIR filter impulse responses, $h_1[n]$, $h_2[n]$, and $h_3[n]$, are plotted in Figure \ref{fig:impulse_responses}.

\begin{figure}[H]
    \centering
    % \includegraphics[width=0.8\textwidth]{figures/impulse_responses.pdf}
    \caption{Impulse responses of the three FIR filters.}
    \label{fig:impulse_responses}
\end{figure}

\subsection{Custom Convolution Implementation (Task 2b)}
To filter the input signal $x[n]$ using the FIR filter $h[n]$, the following convolution operation was implemented:
\begin{equation}
    y[n] = \sum_{k=-\infty}^{\infty} x[k]h[n-k]
\end{equation}
The custom routine includes the \texttt{ylen\_choice} parameter. A verification of this implementation against a toy signal (compared to the standard built-in function) is shown in Figure \ref{fig:convolution_verification}.

\begin{figure}[H]
    \centering
    % \includegraphics[width=0.8\textwidth]{figures/convolution_verification.pdf}
    \caption{Verification of the custom convolution function using a toy signal.}
    \label{fig:convolution_verification}
\end{figure}

\subsection{Discrete-Time Fourier Transform Implementation (Task 2c)}
The Discrete-Time Fourier Transform (DTFT) of the sequences was computed using the definition:
\begin{equation}
    H(e^{j\omega}) = \sum_{n=-\infty}^{\infty} h[n]e^{-j\omega n}
\end{equation}
The implemented function takes the inputs \texttt{h}, \texttt{N}, and \texttt{fs}, and outputs both the complex vector $H(e^{j\omega})$ and the frequency axis in Hertz.

\subsection{Frequency Spectra of Filters (Task 2d)}
The absolute values of the frequency spectra, $|H_1(e^{j\omega})|$, $|H_2(e^{j\omega})|$, and $|H_3(e^{j\omega})|$, are plotted against the physical frequency $F$ in Figure \ref{fig:frequency_spectra}.

\begin{figure}[H]
    \centering
    % \includegraphics[width=0.8\textwidth]{figures/frequency_spectra.pdf}
    \caption{Absolute frequency spectra of the three provided FIR filters.}
    \label{fig:frequency_spectra}
\end{figure}

\subsection{Filter Classification (Task 2e)}
Based on the frequency spectra observed in Figure \ref{fig:frequency_spectra}, the filters are classified as follows:
\begin{itemize}
    \item $h_1[n]$ is a [Insert classification: lowpass/bandpass/highpass] filter.
    \item $h_2[n]$ is a [Insert classification: lowpass/bandpass/highpass] filter.
    \item $h_3[n]$ is a [Insert classification: lowpass/bandpass/highpass] filter.
\end{itemize}

\subsection{Global Data Filtering (Task 2f)}
All 201 input time series were successfully filtered using the three respective impulse responses via built-in convolution routines for execution speed.

\subsection{Section Plot of Propagating Waves (Task 2g)}
To recreate findings similar to Matoza et al. (2022), the data was bandpass filtered between 0.01 and 2 Hz. 
Filter [Insert Filter Name, e.g., $h_2[n]$] was chosen for this task because [Insert your justification here based on your DTFT analysis]. 

Figure \ref{fig:section_plot} shows the resulting section plot with time on the x-axis and distance from Hunga Tonga on the y-axis.

\begin{figure}[H]
    \centering
    % \includegraphics[width=1\textwidth]{figures/section_plot.pdf}
    \caption{Section plot of the filtered waveforms (0.01–2 Hz) demonstrating global wave propagation.}
    \label{fig:section_plot}
\end{figure}

% ---------------------------------------------------------
\section{Task 3: Celerity Estimation for the Acoustic Wave}
% ---------------------------------------------------------

\subsection{Arrival Time Picking (Task 3a)}
A custom routine was developed to loop over all station signals. For each station, the raw and filtered traces were analyzed to manually pick the arrival time of the higher-frequency infrasound wave. The picked times were saved for subsequent analysis.

\subsection{Celerity Estimation and Literature Comparison (Task 3b)}
Figure \ref{fig:celerity_fit} plots the estimated travel times against the event-to-station distances. Identified outliers (caused by noise or competing arrivals) are marked in [Insert Color, e.g., red] and were excluded from the linear regression.

\begin{figure}[H]
    \centering
    % \includegraphics[width=0.8\textwidth]{figures/celerity_fit.pdf}
    \caption{Linear fit of travel time versus distance for celerity estimation. Outliers are excluded from the fit.}
    \label{fig:celerity_fit}
\end{figure}

By fitting a straight line to the accepted distance/time data, the slope yielded an estimated celerity of \textbf{[Insert Value]} m/s. 

Comparing this result to the literature, this estimate aligns with typical celerities of infrasound waves as reported by Vergoz et al. (2022) and Matoza et al. (2022), which state [Insert brief comparison/typical values from the papers].

% ---------------------------------------------------------
\section*{Appendix: Code Links \& Structure}
\addcontentsline{toc}{section}{Appendix: Code Links \& Structure}
% ---------------------------------------------------------
As per the submission requirements, all code is properly commented and attached in the submission archive. 

\begin{itemize}
    \item \textbf{Main Run Script:} \texttt{main.py} / \texttt{main.m} (Running this script without inputs generates all figures included in this report).
    \item \textbf{Filtering Module:} \texttt{filters.py} / \texttt{filters.m} (Contains the custom convolution and DTFT functions).
    \item \textbf{Analysis Module:} \texttt{celerity.py} / \texttt{celerity.m} (Contains the arrival picking tool and linear fit functions).
\end{itemize}

\end{document}